\chapter{Optical Measurement and Test Reference}
\label{app:measurement-reference}

This appendix is a lookup reference for measurement selection, metric
conditions, instruments, stressed-receiver methods, link-budget accounting,
and CMIS diagnostics. It does not define the product-readiness lifecycle;
that story lives in \Cref{ch:product-readiness}.

\textit{Read first:} measurement selection by question; metric and plane table;
TDECQ and SECQ distinctions; consistent link-budget accounting.

\textit{Reference:} instrument quick reference; CMIS state and diagnostics;
Rapid Interview Checks.

\section{Measurement selection by engineering question}
\label{sec:instruments}

Name the question and reference plane before selecting the instrument. Use the
least complex measurement that can falsify the leading hypothesis or support
the required decision.

\begin{description}
\item[Where is the power?] Power meter for average power; DCA for OMA. Walk
      planes before you change bias.
\item[Did modulation quality move?] DCA for PAM4 eyes, TDECQ, OMA, RLM
      (\Cref{sec:tdecq}).
\item[Did the spectrum or grid move?] OSA or wavemeter for wavelength, SMSR,
      and side modes (\Cref{ch:lasers}).
\item[Did receiver margin shift or floor?] BERT with calibrated attenuator or
      stressor; FEC counters for shape
      (\Cref{sec:sensitivity,sec:secq,sec:fec-symbol-histogram}).
\item[Is intensity noise the floor?] PD plus ESA or a dedicated RIN analyzer
      under a defined ORL (\Cref{sec:rin-values}).
\item[Is the error random, bursty, or state-driven?] Pre-FEC BER plus FEC
      histograms and management state (\Cref{sec:fec-symbol-histogram,sec:cmis}).
\item[Does management telemetry match bench truth?] Host or CMIS tool versus
      external meter at a named plane (\Cref{sec:cmis}).
\end{description}

Use electrical loopback, optical loopback, and golden-host or golden-module
swaps to bisect ownership before opening supplier FA.

\section{Metric and reference-plane table}
\label{sec:measurement-mapping}

\begin{table*}[ht]
\refstepcounter{table}\label{tab:measurement-mapping}%
\centering
\setlength{\tabcolsep}{3pt}%
\footnotesize
\renewcommand{\arraystretch}{1.2}%
\begin{tabular}{@{}p{0.16\linewidth}p{0.22\linewidth}p{0.18\linewidth}p{0.20\linewidth}p{0.18\linewidth}@{}}
\toprule
Metric & Required condition and plane & Typical class & Establishes & Does not establish \\
\midrule
OMA / TDECQ & Named pattern, EQ, temperature; Tx optical plane & Tx quality &
Modulated quality vs reference receiver & Unique physical mechanism \\
\addlinespace[0.6em]
ER / RLM & Same plane as OMA & Tx levels & Level structure and usable swing &
Receiver sensitivity alone \\
\addlinespace[0.6em]
Wavelength / SMSR & Grid window; spectral plane & Spectral & Grid placement and
mode purity & BER floor ownership \\
\addlinespace[0.6em]
RIN & Stated ORL and bias quietness & Intensity noise & Signal-proportional
noise contribution & That every BER floor is RIN \\
\addlinespace[0.6em]
IL / ORL & Connector class, plant, cleanliness & Channel & Power and reflection
assumptions & Signal-quality closure \\
\addlinespace[0.6em]
Rx sensitivity & Plane, BER/FEC, pattern, stress & Receiver & Minimum OMA for
objective & Unstressed field plant alone \\
\addlinespace[0.6em]
Pre-FEC BER / FEC histogram & Named FEC and dwell & Link/system & Error rate and
burst vs sparse shape & Confirmed mechanism alone \\
\addlinespace[0.6em]
CMIS state / DDM & Host sequence; named monitors & Control & Management and
diagnostic path & Optical mechanism by itself \\
\bottomrule
\end{tabular}
\par\medskip
{\raggedright\footnotesize\textbf{Table~\thetable.} Metric lookup: conditions,
what the measurement establishes, and what it does not. Instrument choices are
in \Cref{sec:instruments}.\par}
\end{table*}

\paragraph{Channel evidence.}
\label{sec:optical-channel}
Insertion loss, ORL, dispersion, and filtering belong with the plant
assumptions used in the link budget. Reflections can seed feedback noise, MPI,
distortion, and power-independent floors even when average power is stable
(\Cref{sec:rin-values,ch:lasers}).

\section{TDECQ reference}
\label{sec:tdecq}

\term{TDECQ} (transmitter and dispersion eye closure quaternary) asks how much
worse a PAM4 transmitter is than an ideal source after a defined reference
receiver and bounded equalizer. Use the exact PMD clause for filter, equalizer
limits, histogram locations, and target error ratio.

Representative procedure: capture the optical waveform on a DCA through the
clause reference receiver; apply the reference FFE; evaluate noise at the
required PAM4 thresholds; report the dB ratio of ideal versus measured
tolerable noise. Lower is better; the numeric cap is PMD-specific.

TDECQ is a composite metric. Uneven levels point toward modulator or driver
linearity (RLM). Residual eye closure the equalizer cannot remove points toward
excess ISI or limited bandwidth. A noise-limited result points toward low OMA,
RIN, or reflections. Passing average power does not establish a valid PAM4
transmitter. Do not double-count TDECQ in a link budget that already embeds the
compliance OMA/TDECQ relationship (\Cref{sec:link-budget}).

\section{SECQ and stressed-receiver reference}
\label{sec:secq}

\term{SECQ} (stressed eye closure quaternary) is a \emph{receiver-side}
stressed-eye method under a named specification. It applies a calibrated optical
stressor and asks how much margin remains before the receiver hits that clause's
target pre-FEC metric.

TDECQ evaluates transmitter quality through a reference-receiver model. SECQ is
not the same test as a generic stressed-receiver sensitivity sweep, even when
both use attenuation and ISI. Always state the PMD, FEC architecture, error
model, metric, stress calibration, and test duration. For LPO, stressed Rx
margin on the host-side receiver is as important as transmitter-quality metrics
(\Cref{sec:equalization,sec:conditioning}).

\section{Link-budget accounting reference}
\label{sec:link-budget}

A link budget is a signed ledger from transmitter to receiver. For IM/DD short
reach, start from outer OMA at a named Tx plane, subtract each loss once, and
compare the remainder to receiver sensitivity at the named BER or FEC objective
(\Cref{sec:kp4,sec:sensitivity}).

\begin{dectree}
Transmitter output (OMA)
  |
Coupling loss
  |
Connector loss
  |
Fiber / waveguide loss
  |
Transmitter-quality and other penalties (one convention)
  |
Receiver input
  |
Sensitivity requirement
  |
Remaining margin
\end{dectree}

Use one internally consistent convention. State whether transmitter-quality
effects are embedded in the compliance requirement or represented as a separate
engineering penalty, and never count the same impairment twice. Do not mix
average-power and OMA budgets, reference planes, or embedded and explicit
transmitter penalties.

Keep power, signal-quality, timing, thermal, and control authority as separate
ledgers when the impairment is not a pure optical-power number
(\Cref{sec:laser-margin-erosion,sec:tree-margin-budget}). Distinguish design
allocation from measured net margin across the operating envelope.

\section{Instrument quick reference}

\begin{description}
\item[Power meter] Average optical power at a named plane; does not measure OMA
      or eye quality.
\item[DCA / sampling oscilloscope] Eyes, TDECQ, OMA, RLM with a reference
      receiver matched to the PHY.
\item[BERT and FEC counters] Pre- and post-FEC BER, dwell, and histogram shape.
\item[VOA / calibrated stressor] Controlled attenuation and optional ISI for
      sensitivity and stressed-receiver work.
\item[OSA / wavemeter] Wavelength, SMSR, side modes.
\item[ORL / reflection setup] Plant return-loss condition for RIN and burst
      investigations.
\item[RIN measurement] PD plus ESA or RIN analyzer under defined ORL and quiet
      bias.
\item[Chamber / TEC] Temperature corners for reversible operation and lock
      behavior.
\item[CMIS / host tools] Management state, DDM correlation, alarms, and
      recovery.
\end{description}

\section{CMIS and diagnostic reference}
\label{sec:cmis}

\term{CMIS}\firstterm{CMIS} (Common Management Interface Specification) is the
vendor-neutral management layer between a host and a module that implements it.
Form factors and optional pages vary; do not imply identical behavior across
pluggables, ELSFP, and CPO engines~\cite{cmis53}.

The host drives a module state machine toward ModuleReady before authorizing
light. Data-path and network-path states in later CMIS revisions refine lane
enable. A module forced into an emitting state that passes BER has not passed
bring-up if the required management sequence, safe state, diagnostics, alarms,
and recovery behavior are incorrect (\Cref{sec:bringup}).

\term{DDM}\firstterm{DDM} provides per-lane Tx/Rx power, bias when exposed,
temperature, rails, LOS/LOL, and alarms. On bring-up, dump the register map you
will use in the field and treat disagreement between DDM and an external meter
as a finding. ATP should prove ModuleReady across voltage and thermal corners
and ECO-control firmware like other production revisions
(\Cref{sec:supplier-exec}).

Operational use of management state is decision-oriented, not a register
inventory (\Cref{ch:networking,sec:networking-ops}):

\begin{itemize}
\item Management state must be trustworthy before counters drive containment.
\item Every counter needs a semantic definition, accumulation window, and reset
      behavior.
\item Polling cadence sets what bursts and short recoveries you can observe.
\item A reset or clear can erase the failing evidence; preserve dumps before
      recovery when the event matters.
\item Alarm latching and clearing must be understood before a ``cleared'' flag
      is treated as healthy.
\end{itemize}

Exact page, bank, and bit maps remain vendor- and revision-specific; keep them
in the product management reference, not in the operational chapter.

\section{Rapid Interview Checks}

\paragraph{Prompt.} What does a passing average-power measurement fail to
establish?\\
\textit{Check.} Modulation quality, eye closure, level linearity, and
signal-quality margin.

\paragraph{Prompt.} What must accompany a sensitivity number?\\
\textit{Check.} Optical reference plane, BER or FEC condition, pattern,
temperature, wavelength, and stress method when claimed
(\Cref{ch:product-readiness,sec:sensitivity}).

\paragraph{Prompt.} How does SECQ differ from generic stressed sensitivity?\\
\textit{Check.} SECQ is a named stressed-eye method under a clause; generic
stressed sensitivity is not automatically SECQ (\Cref{sec:secq}).

\paragraph{Prompt.} What does a locked flag fail to prove?\\
\textit{Check.} Margin, interoperability, and that the supported plant and host
corners remain closed.

\paragraph{Prompt.} Why can TDECQ be double counted?\\
\textit{Check.} Embedding compliance OMA/TDECQ and also subtracting an
independent TDECQ penalty taxes the same impairment twice
(\Cref{sec:link-budget}).

\paragraph{Prompt.} Which decision should precede instrument selection?\\
\textit{Check.} Name the question, population, condition, and reference plane;
then choose the least complex measurement that changes the decision
(\Cref{sec:instruments}).
